\documentclass[french]{article}
\usepackage{verbatim}
\usepackage{amsmath,amsfonts,amssymb}
\usepackage{pgfplots}
\usepackage[utf8]{inputenc}
\usepackage[french]{babel}
\usepackage{pgf,tikz,tkz-tab}
\usepackage{mathrsfs}
\usetikzlibrary{arrows}
\usepackage{pstricks}
\pgfplotsset{compat=newest}
\usepackage{geometry}
\geometry{hmargin=1cm,vmargin=2cm}


\begin{document}
\begin{center}
\underline{\textbf{Un problème de l'université d'Oxford}}
\end{center}

\section{Un stylo qui glisse}

On prend un stylo posé à la verticale contre un mur, ce stylo glisse contre le mur jusqu'à être sur le sol. On s'intéresse à la trajectoire du milieu du stylo que l'on représentera par un point noté M  qui aura pour coordonnée $(x;y)$.
On suppose que le stylo est de longueur $L$.  

\definecolor{qqwuqq}{rgb}{0,0.39215686274509803,0}
\definecolor{xdxdff}{rgb}{0.49019607843137253,0.49019607843137253,1}
\definecolor{uuuuuu}{rgb}{0.26666666666666666,0.26666666666666666,0.26666666666666666}
\begin{tikzpicture}[scale=1,line cap=round,line join=round,>=triangle 45,x=1cm,y=1cm]
\begin{axis}[
x=1cm,y=1cm,
axis lines=middle,
xmin=-5.866379805052241,
xmax=13.547501529625054,
ymin=-2.956904761844938,
ymax=13.301132633899238,
xtick={-4,-2,...,12},
ytick={-2,0,...,12},]
\clip(-5.866379805052241,-2.956904761844938) rectangle (13.547501529625054,13.301132633899238);
\draw[line width=2pt,color=qqwuqq,fill=qqwuqq,fill opacity=0.10000000149011612] (3.437229331043156,0) -- (3.437229331043156,0.4616935137074414) -- (2.9755358173357145,0.4616935137074414) -- (2.9755358173357145,0) -- cycle; 
\draw [line width=2pt] (6,0)-- (-0.04892836532857103,10.363902046580883);
\draw [line width=2pt,dash pattern=on 1pt off 1pt on 1pt off 4pt] (2.9755358173357145,5.181951023290441)-- (0,5.181951023290441);
\draw [line width=2pt,dash pattern=on 1pt off 1pt on 1pt off 4pt] (2.9755358173357145,5.181951023290441)-- (2.9755358173357145,0);
\begin{scriptsize}
\draw [fill=uuuuuu] (6,0) circle (2pt);
\draw[color=uuuuuu] (6.169356044809355,0.42746580714866256) node {$B$};
\draw [fill=xdxdff] (-0.04892836532857103,10.363902046580883) circle (2.5pt);
\draw[color=xdxdff] (0.11884145844131944,10.830868585148284) node {$A$};
\draw [fill=uuuuuu] (2.9755358173357145,5.181951023290441) circle (2pt);
\draw[color=uuuuuu] (3.1440987516253376,5.607402755190314) node {$M$};
\draw [fill=uuuuuu] (4.487767908667857,2.5909755116452207) circle (2pt);
\end{scriptsize}
\end{axis}
\end{tikzpicture}

\begin{enumerate}
\item Compléter les longueurs sur la figure.
\item À l'aide du théorème de Pythagore, en déduire une équation entre $x$,$y$ et $L$.
\item On admet que l'équation cartésienne d'un cercle de centre O et de rayon $r$ est $x^2+y^2=r^2$, quel est la trajectoire du point M ? 
\end{enumerate}
\newpage
\section{Un stylo qui tombe}

On prend un stylo posé à la verticale contre un mur, ce stylo tombe du mur jusqu'à être sur le sol. On s'intéresse à la trajectoire du milieu du stylo que l'on représentera par un point noté M  qui aura pour coordonnée $(x;y)$.
On suppose que le stylo est de longueur $L$.  

\definecolor{qqwuqq}{rgb}{0,0.39215686274509803,0}
\definecolor{xdxdff}{rgb}{0.49019607843137253,0.49019607843137253,1}
\definecolor{uuuuuu}{rgb}{0.26666666666666666,0.26666666666666666,0.26666666666666666}
\begin{tikzpicture}[line cap=round,line join=round,>=triangle 45,x=1cm,y=1cm]
\begin{axis}[
x=1cm,y=1cm,
axis lines=middle,
xmin=-1.111718637385122,
xmax=17.160076230940884,
ymin=-2.25503775190276,
ymax=13.046566246571869,
xtick={0,2,...,16},
ytick={-2,0,...,12},]
\clip(-1.111718637385122,-2.25503775190276) rectangle (17.160076230940884,13.046566246571869);
\draw[line width=2pt,color=qqwuqq,fill=qqwuqq,fill opacity=0.10000000149011612] (4.059422363593717,0.434532849411756) -- (3.6248895141819615,0.4345328494117561) -- (3.6248895141819615,0) -- (4.059422363593717,0) -- cycle; 
\draw [line width=2pt] (0,0)-- (8.118844727187435,8.836535536952296);
\draw [line width=2pt,dash pattern=on 1pt off 1pt on 1pt off 4pt] (4.059422363593717,4.418267768476148)-- (4.059422363593717,0);
\begin{scriptsize}
\draw [fill=uuuuuu] (0,0) circle (2pt);
\draw[color=uuuuuu] (0.15829401041332236,0.39764995599879893) node {$B$};
\draw [fill=xdxdff] (8.118844727187435,8.836535536952296) circle (2.5pt);
\draw[color=xdxdff] (8.290471771316264,9.267254415623317) node {$A$};
\draw [fill=uuuuuu] (4.059422363593717,4.418267768476148) circle (2pt);
\draw[color=uuuuuu] (3.6405867543767982,4.904146448186775) node {$M$};
\draw [fill=uuuuuu] (4.059422363593717,0) circle (2pt);
\end{scriptsize}
\end{axis}
\end{tikzpicture}

\begin{enumerate}
\item Compléter les longueurs sur la figure.
\item À l'aide du théorème de Pythagore, en déduire une équation entre $x$,$y$ et $L$.
\item On admet que l'équation cartésienne d'un cercle de centre O et de rayon $r$ est $x^2+y^2=r^2$, quel est la trajectoire du point M ? 
\item Que remarque t-on entre le milieu du stylo qui tombe et le milieu du stylo qui glisse?
\end{enumerate}

\section{Simulation dynamique}
Pour une animation, rendez-vous sur:
https://www.desmos.com/calculator/86hptwtqwb

\end{document}
